\chapter{INTRODUCTION}

\section{BACKGROUND}

A battery management system (BMS) protects a lithium-ion pack and estimates its state of charge. It already holds the three measurements from which ageing is visible: terminal voltage, current and temperature \cite{waag2014}. A software layer above it can, in principle, turn those measurements into answers that an operator or driver can act on: how much of its original capacity the battery retains (state of health, SOH), how fast that is changing, how long it has before it should be replaced (remaining useful life, RUL), and what usage is doing to it.

Capacity and resistance are the two quantities through which ageing usually appears. Capacity fades as lithium inventory and active material are lost, and internal resistance grows, which reduces the power the cell can deliver \cite{vetter2005,barre2013}. Reviews of SOH estimation distinguish direct measurement, model-based estimation and data-driven approaches \cite{berecibar2016,waag2014}. A vehicle, however, rarely provides the conditions a laboratory does: discharges are partial, rates vary, and no reference test is available. Any practical health estimator must therefore work from what the vehicle actually records.

\section{MOTIVATION}

The dominant way to build such a software layer is to fit a model. Features are extracted from cycling data and a regressor is trained to predict SOH, with feature sets ranging from usage aggregates to discharge-curve descriptors such as the capacity--voltage difference between cycles \cite{severson2019}. Skill is commonly reported under leave-one-cell-out validation, in which a held-out cell's siblings from the same test protocol remain in the training set, so the test asks whether the model recognises a protocol it has already seen.

A deployed BMS faces a harder question: the next battery comes from conditions the training set never contained. This project began as a battery-degradation study built on fitted models, and measured that question directly. On a benchmark of 31 NASA cells in nine protocol cohorts, eleven of twelve methods lost skill, and the twelfth, a mean-only baseline, could not when the held-out unit changed from a cell to a whole protocol cohort, and no method could be distinguished from another (Chapter~\ref{ch:results}). That negative result is the motivation for what follows: if models fitted across cells do not transfer, a health layer for a vehicle has to be built from estimators that fit nothing across cells, and from a pipeline that is honest about when it does not know.

The work is also motivated by engineering practice. A monitoring system that is asked for a number it cannot compute and returns a plausible one is dangerous. A defect of exactly this kind occurred during development: a missing temperature channel was scored as ``temperature is fine'', because a comparison against \texttt{NaN} evaluates to false, so every heat check silently passed for a pack nobody had measured. The design of BEACON takes the opposite stance: the system refuses to produce a number it cannot justify, and says why.

\section{PROBLEM STATEMENT}

The problem addressed in this project is how to build a software layer above a BMS that can measure battery health from the BMS's own signals, namely voltage, current and time, without depending on a model trained on other cells, and that knows when its evidence is insufficient. Three sub-problems follow.

\begin{itemize}
    \item \textbf{Measurement.} How can SOH, internal resistance and RUL be measured from voltage, current and time alone, including from partial discharges, and how accurately?
    \item \textbf{Honesty.} How can the system decide, without ground truth, when an estimate should be reported with a stated error band and when it should be refused?
    \item \textbf{Integration.} How can data from a CAN bus, a serial bench rig and research datasets all be scored through one shared pipeline and served to a dashboard, so that a difference in output is a difference in the telemetry and not in the code?
\end{itemize}

\section{RESEARCH GAP}

Data-driven SOH estimators are widespread, but how they are validated changes what they appear to achieve. Group-wise splitting by cell changes reported skill substantially \cite{maher2025}, and a large gap between random and cell-wise splits has been reported on NASA data \cite{le2026}. These studies stop at the cell. What remains insufficiently characterised is the next step: whether skill survives a change of \emph{test protocol}, and what a BMS software layer can still measure reliably when it does not.

Curve-based methods, including differential voltage analysis \cite{bloom2005} and incremental capacity analysis, including on-board use \cite{weng2013}, read degradation from the shape of the voltage--charge curve. Online recursive estimation of cell parameters, including resistance, is established \cite{plett2004}. The gap this project addresses is the combination: estimators validated from raw telemetry rather than from laboratory summaries, with an online resistance correction taken from the load step, gates that refuse rather than extrapolate, and a measured criterion for when the evidence suffices, all evaluated on more than one dataset without re-tuning.

\section{OBJECTIVES}

The objectives of the project are:

\begin{enumerate}
    \item To build an end-to-end telemetry system that ingests data from a CAN bus, a serial bench rig and research datasets, converges them on one schema, and scores them through one shared pipeline.
    \item To estimate SOH from voltage, current and time alone, with an ohmic correction identified from the load step, and to validate it against laboratory capacity.
    \item To report resistance as a second validated health quantity, and to estimate RUL from each cell's own fade trend, reported as a user meets it.
    \item To measure rather than assert when the evidence suffices, so that every output is validated, banded by its measured error, or refused.
    \item To validate externally, without re-tuning, on a second and third dataset, and to retain negative results.
    \item To build and bring up a physical data-acquisition rig, and to serve all results through a REST API and a web dashboard.
\end{enumerate}

\section{BATTERY HEALTH INDICATORS}

Three quantities are used throughout this report, and they are defined here once so that every later chapter uses them in the same way.

\subsection{State of Health}

State of health is the fraction of the original capacity that the cell can still deliver. With $Q_{\mathrm{now}}$ the capacity measured on a full constant-current discharge and $Q_{\mathrm{ref}}$ the reference capacity of the same cell, it is
\begin{equation}
	\mathrm{SOH} = \frac{Q_{\mathrm{now}}}{Q_{\mathrm{ref}}} \times 100\,\%.
	\label{eq:sohdef}
\end{equation}
The choice of $Q_{\mathrm{ref}}$ matters. It may be the rated capacity from the data sheet or the capacity measured on the cell's own first discharges. The two give different numbers for the same cell, and this report states which is used wherever a figure is quoted. A battery is conventionally regarded as having reached end of life when SOH falls to 80\%. The RUL estimator of this project uses a default threshold of 90\%, because the CALCE trajectories rarely reach 80\% within the recorded life, and reports the 80\% case separately where the data allow it.

\subsection{Internal Resistance}

Internal resistance is the ratio of the voltage step to the current step when a load is applied. For a step from current $I_{1}$ to $I_{2}$ that moves the terminal voltage from $V_{1}$ to $V_{2}$,
\begin{equation}
	R_{\mathrm{int}} = \frac{\left|V_{2}-V_{1}\right|}{\left|I_{2}-I_{1}\right|}.
	\label{eq:rintdef}
\end{equation}
Resistance rises as the cell ages, and it determines the voltage sag under load and therefore the power the cell can deliver.

\subsection{Remaining Useful Life}

Remaining useful life is the number of further cycles before SOH reaches the end-of-life threshold. It is a prediction, so it carries an error that grows with the distance from the threshold, and a figure for it is meaningful only if the error is stated alongside it.

\section{PROPOSED APPROACH IN BRIEF}

BEACON, the system built in this project, is organised in three tiers. The data tier receives telemetry from three sources, a controller area network (CAN) bus, a serial link from a bench rig and public research datasets, and converts all three into one unified frame. The analytics tier is a Python service that scores every frame through a single function and returns state of health, resistance, remaining useful life and a behaviour risk score, each with its evidence and, where the evidence is insufficient, a refusal with the reason. The presentation tier is an Express gateway and a React dashboard that display the results and label every value with its source.

Three design commitments run through the work. The first is that nothing is fitted across cells: the estimators measure each cell against its own early behaviour. The second is that a result is either validated against laboratory truth, reported with a measured error band, or refused. The third is that the same scoring function serves a vehicle, a rig and a dataset, so that any difference between their outputs comes from the telemetry and not from the code.

\section{SIGNIFICANCE OF THE WORK}

For a fleet operator or an owner, a health figure is useful only if it can be trusted, and a figure that cannot be trusted is worse than none. By measuring its own error on public data and by declining to answer when its evidence is poor, the system gives an operator a reason to act on an estimate when one is shown and a reason to collect more data when one is withheld. For the research community, the retained negative results, in particular the loss of skill of fitted models when the test protocol changes, are a reusable finding about how battery health models should be validated. For students and practitioners, the repository is a worked example of how a software layer over a BMS can be tested, documented and reproduced.

\section{SCOPE AND ORGANISATION OF THE REPORT}

Development and primary validation use single lithium cobalt oxide (LCO) cells at constant current (CALCE). External validation uses NMC/LCO blend pouch cells at 40~$^\circ$C (Oxford), which also supply one real drive-cycle discharge. A third dataset (NASA) is used for the benchmark and for a cross-dataset test. The project makes no claim about packs, lithium iron phosphate, or temperature-dependent capacity correction. The bench rig validates data acquisition, not health estimation: no discharge has yet been recorded on it.

Table~\ref{tab:chapters} gives a map of the chapters. The remainder of the report is organised as follows. Chapter~\ref{ch:lit} reviews related work. Chapter~\ref{ch:method} presents the system design and the estimation methods. Chapter~\ref{ch:implementation} describes the implementation, the hardware rig and the deployment. Chapter~\ref{ch:results} reports the results and discusses them. Chapter~\ref{ch:conclusion} concludes and lists future work.

\begin{table}[htbp]
	\caption{Organisation of the Report}
	\centering
	\small
	\renewcommand{\arraystretch}{1.15}
	\setlength{\tabcolsep}{4pt}
	\begin{tabular}{|C{0.1\linewidth}|L{0.28\linewidth}|L{0.44\linewidth}|}
		\hline
		\textbf{Chapter} & \textbf{Title} & \textbf{Content} \\
		\hline
		1 & Introduction & Background, problem, objectives and scope \\
		\hline
		2 & Literature Review & Ageing, estimation methods, interfaces and validation practice \\
		\hline
		3 & Methodology / System Design & Requirements, schema, estimators, gates and validation protocol \\
		\hline
		4 & Implementation and Deployment & Acquisition, rig, services, security and testing \\
		\hline
		5 & Results and Discussion & Benchmark, estimator accuracy, hardware results and threats to validity \\
		\hline
		6 & Conclusion and Future Work & Findings, limitations, lessons and next steps \\
		\hline
	\end{tabular}
	\label{tab:chapters}
\end{table}
